\chapter{SHBG (Sex Hormone Binding Globulin)}

\section*{What Is This Test?}
Sex hormone-binding globulin (SHBG) is a 373-amino-acid glycoprotein produced primarily by the liver. SHBG binds with high affinity to testosterone and dihydrotestosterone (DHT), and with approximately 2--3 times lower affinity to estradiol. The principal function of SHBG is to transport sex steroids in the blood and regulate their bioavailability. In circulation, approximately 44--60\% of testosterone is tightly bound to SHBG and essentially unavailable for biological action; 38--54\% is weakly bound to albumin and can dissociate; and 1--2\% is free (unbound). Only the free and albumin-bound fractions (together called bioavailable testosterone) can enter target cells and exert biological effects. SHBG levels are regulated by multiple hormones and metabolic factors: estrogens, thyroid hormone, and liver disease increase SHBG production; androgens, insulin, insulin-like growth factor 1 (IGF-1), and obesity decrease SHBG production. SHBG concentrations vary 10-fold in health and disease. Measurement of SHBG is clinically essential because total testosterone concentrations are largely determined by SHBG levels, and cannot be reliably interpreted without knowledge of SHBG. The SHBG test is used in conjunction with total testosterone to calculate free and bioavailable testosterone, to evaluate hyperandrogenism in PCOS, to explain discrepancies between total testosterone and clinical androgen status, and as a marker of insulin resistance.

\section*{Alternative Names}
Sex Steroid-Binding Globulin; Sex Hormone-Binding Protein; Testosterone-Estradiol Binding Globulin (TeBG, historical); SHBG, Serum; Free Androgen Index (FAI = testosterone/SHBG x 100)

\section*{Principle}
SHBG is measured using a two-site sandwich immunometric (chemiluminescent) immunoassay on automated analyzers. Two monoclonal antibodies against different epitopes on the SHBG molecule are used: a capture antibody on a solid phase and a detection antibody labeled with a chemiluminescent tag (ruthenium on Roche ECLIA, acridinium ester on Abbott/Siemens). The signal is proportional to SHBG concentration. Common platforms: Roche Cobas ECLIA and Abbott Architect CMIA. SHBG assays have good precision (CV <5\%) and are calibrated against the WHO 2nd International Standard (IS 08/266). SHBG immunoassays measure total immunoreactive SHBG protein and do not distinguish between occupied (steroid-bound) and unoccupied SHBG. SHBG is relatively stable and does not require special pre-analytical handling.

\section*{History}
SHBG was first identified in the 1960s as a plasma protein that bound testosterone with high affinity, distinct from albumin. It was initially called testosterone-estradiol binding globulin (TeBG) before being renamed SHBG in the 1970s to reflect its broader steroid-binding properties. The radioimmunoassay for SHBG was developed in the 1980s. The clinical importance of SHBG in interpreting testosterone levels was established by Vermeulen, Rosner, and others in the 1990s, who developed equations to calculate free testosterone from total testosterone, SHBG, and albumin. The role of low SHBG as a marker of insulin resistance and metabolic syndrome was recognized in the 1990s--2000s, and SHBG is now considered a surrogate marker for insulin resistance in PCOS. The Free Androgen Index (FAI = total testosterone/SHBG x 100) was introduced as a simple measure of androgen excess that does not require albumin measurement.

\section*{Why This Test Is Ordered}
\begin{itemize}
  \item Calculation of free and bioavailable testosterone from total testosterone in men being evaluated for hypogonadism, especially when total testosterone is borderline (250--350 ng/dL) or when conditions altering SHBG are present (obesity, diabetes, aging, liver disease, HIV, medications)
  \item Evaluation of hyperandrogenism in women with PCOS, hirsutism, or oligomenorrhea --- low SHBG is a characteristic finding in PCOS due to hyperinsulinemia suppressing hepatic SHBG production
  \item Calculation of the free androgen index (FAI) as a measure of androgen excess --- FAI >4--5 in women suggests hyperandrogenism
  \item Investigation of discordance between total testosterone and clinical signs of androgen deficiency or excess
  \item Assessment of insulin resistance --- low SHBG independently predicts insulin resistance, metabolic syndrome, and type 2 diabetes risk, even in individuals with normal total testosterone
  \item Evaluation of precocious puberty or delayed puberty --- SHBG levels decrease with pubertal onset due to rising androgens and insulin
  \item Monitoring of conditions that profoundly alter SHBG: estrogen therapy (increases SHBG 2--4 fold), oral contraceptives, thyroid dysfunction, liver disease, androgen therapy (decreases SHBG), glucocorticoid therapy
  \item Investigation of gynecomastia in males --- SHBG abnormalities may contribute to an altered estradiol-to-testosterone ratio
  \item Evaluation of women with clinical hyperandrogenism (hirsutism, acne, androgenic alopecia) whose total testosterone is normal --- normal total testosterone with low SHBG may still result in elevated free testosterone
\end{itemize}

\section*{Primary vs Secondary Test}
SHBG is a secondary (adjunctive) test. It is never interpreted in isolation but is used to interpret total testosterone and calculate free testosterone. In PCOS evaluation, low SHBG is a supportive finding and a marker of insulin resistance. SHBG is an essential component of the evaluation of androgen disorders because total testosterone alone cannot be interpreted without knowledge of SHBG binding capacity.

\section*{How This Test Is Performed}
A healthcare professional draws a blood sample from a vein into a serum separator tube (gold-top) or plain red-top tube (~3--5 mL). SHBG does not have significant circadian variation (unlike testosterone), so the draw can be at any time of day, though morning collection alongside testosterone (which does have circadian variation) is standard. SHBG is stable at room temperature for 24 hours and refrigerated for 7 days.

\section*{What Preparation Is Needed}
No fasting is required. However, SHBG should be measured simultaneously with total testosterone for calculation of free testosterone. High-dose biotin (>5 mg/day) must be held 48 hours before testing. Estrogen therapy and oral contraceptives markedly increase SHBG (2--4 fold); these medications should be documented and, if possible, held for 4--6 weeks before testing for evaluation of androgen status. Androgen therapy decreases SHBG. Insulin, IGF-1, glucocorticoids, and progestins with androgenic activity decrease SHBG. Thyroid hormone increases SHBG; untreated hypothyroidism decreases SHBG. Severe liver disease decreases SHBG production.

\section*{Contraindications and Precautions}
SHBG levels are profoundly affected by multiple endocrine, metabolic, and pharmacologic factors. Conditions that increase SHBG: estrogen therapy (oral contraceptives, HRT --- 2--4 fold increase), pregnancy (5--10 fold increase by third trimester), hyperthyroidism, liver cirrhosis, hepatitis, HIV infection (possibly related to cytokines), anticonvulsants (phenytoin, carbamazepine, phenobarbital), tamoxifen, and aging. Conditions that decrease SHBG: obesity (inversely proportional to BMI), type 2 diabetes, insulin resistance/metabolic syndrome, PCOS, exogenous androgens/anabolic steroids, hypothyroidism, glucocorticoid excess (Cushing syndrome, pharmacologic), nephrotic syndrome (urinary loss of SHBG), acromegaly (GH/IGF-1 suppression of SHBG), and progestins with androgenic activity (norethindrone, levonorgestrel). SHBG levels are approximately 30--50\% higher in females than males due to estrogen stimulation, and about 2-fold higher in children than adults (declining at puberty with rising androgens). In pregnancy, SHBG increases 5--10 fold by the third trimester, which is a normal physiologic adaptation to protect the mother and fetus from the very high circulating androgens and estrogens. Racial differences: SHBG is higher in Asian individuals compared to Caucasians and African Americans, which should be considered when using reference ranges. Free androgen index (FAI) = total testosterone (nmol/L) / SHBG (nmol/L) x 100. FAI >4--5 in women suggests hyperandrogenism. FAI should NOT be used in men because the relationship between testosterone and SHBG becomes non-linear at male testosterone concentrations.

\section*{Risks}
Minimal risk. Slight pain, bruising, or hematoma at venipuncture site. Rarely: vasovagal syncope.

\section*{What Happens After the Test}
Results available within 1--4 hours. SHBG is interpreted together with total testosterone to calculate free testosterone (Vermeulen equation: using total testosterone, SHBG, and albumin) or free androgen index. In men with borderline total testosterone (250--350 ng/dL): if SHBG is elevated (e.g., >50 nmol/L), free testosterone is likely low and the patient has true hypogonadism. If SHBG is low (e.g., <15 nmol/L), free testosterone is likely normal and the low total testosterone is due to low SHBG rather than hypogonadism. In women with PCOS: SHBG is typically low (<30 nmol/L) due to hyperinsulinemia. Low SHBG with normal total testosterone still results in elevated free testosterone and clinical hyperandrogenism. In patients with very high SHBG (estrogen therapy, pregnancy, hyperthyroidism): total testosterone is elevated but free testosterone is normal; treatment for androgen excess is not needed.

\section*{Understanding Results}

\subsection*{Below Normal}
\begin{itemize}
  \item Obesity: SHBG is inversely proportional to BMI. Adipose tissue produces inflammatory cytokines (TNF-alpha, IL-1) that suppress hepatic SHBG synthesis. Insulin directly suppresses SHBG gene transcription. Weight loss increases SHBG.
  \item PCOS and insulin resistance: low SHBG is a hallmark of PCOS, correlating with the degree of insulin resistance and hyperandrogenism. Low SHBG contributes to elevated free testosterone even when total testosterone is normal.
  \item Type 2 diabetes and metabolic syndrome: low SHBG independently predicts the development of type 2 diabetes, with a stronger predictive value than total testosterone. Low SHBG is associated with all components of the metabolic syndrome.
  \item Hypothyroidism: decreased hepatic SHBG synthesis. Reverses with levothyroxine therapy.
  \item Exogenous androgens/anabolic steroids: suppress SHBG production. Testosterone therapy in hypogonadal men lowers SHBG by 20--30\%.
  \item Glucocorticoid excess: Cushing syndrome or pharmacologic glucocorticoid therapy suppresses SHBG.
  \item Acromegaly: elevated GH and IGF-1 suppress SHBG production.
  \item Nephrotic syndrome: urinary loss of SHBG protein.
  \item Non-alcoholic fatty liver disease (NAFLD): low SHBG is associated with hepatic steatosis and may be a causative factor (SHBG protects against hepatic lipogenesis).
\end{itemize}

\subsection*{Above Normal}
\begin{itemize}
  \item Estrogen therapy and oral contraceptives: estrogen stimulates hepatic SHBG production 2--4 fold. This is the most common cause of elevated SHBG in clinical practice. Oral estrogen has a greater effect than transdermal estrogen because of the first-pass hepatic effect.
  \item Pregnancy: SHBG increases 5--10 fold by the third trimester as a protective mechanism against maternal and fetal hyperandrogenemia.
  \item Hyperthyroidism: thyroid hormone directly stimulates SHBG gene transcription. SHBG normalizes with anti-thyroid therapy.
  \item Liver cirrhosis and hepatitis: paradoxically increase SHBG despite decreased hepatic synthetic function, possibly due to estrogen excess (impaired hepatic estrogen metabolism) or altered cytokine milieu.
  \item Aging in men: SHBG increases progressively with age (1--2\% per year), while total testosterone declines. The combination results in a more rapid decline in free testosterone than total testosterone.
  \item HIV infection: SHBG is elevated, possibly due to altered cytokine production, despite low testosterone.
  \item Anticonvulsant medications: phenytoin, carbamazepine, and phenobarbital induce hepatic SHBG synthesis.
  \item Tamoxifen: selective estrogen receptor modulator with estrogenic activity in the liver, increasing SHBG.
  \item Constitutional delay of puberty: SHBG is normally high in pre-pubertal children and declines at puberty. Failure to decline suggests hypogonadism.
  \item Anorexia nervosa and severe malnutrition: SHBG is elevated due to low insulin and IGF-1, despite low estrogen.
  \item Idiopathic elevation: some individuals have constitutionally high SHBG, which can cause elevated total testosterone without hyperandrogenism.
\end{itemize}

\section*{Normal Results}
\begin{itemize}
  \item SHBG, adult male (20--49 years): 10--57 nmol/L (0.3--1.6 mcg/mL)
  \item SHBG, adult male (>50 years): 20--70 nmol/L (age-related increase)
  \item SHBG, adult female (premenopausal, non-pregnant): 18--144 nmol/L (0.5--4.1 mcg/mL); typical range 30--90 nmol/L in regularly cycling women
  \item SHBG, adult female (postmenopausal): 20--100 nmol/L (slightly lower than premenopausal)
  \item SHBG, pregnancy (third trimester): 200--500 nmol/L (5--10 fold increase)
  \item SHBG, children (pre-pubertal): 50--150 nmol/L (higher than adults)
  \item SHBG in PCOS: typically <30 nmol/L (low)
  \item Free androgen index (FAI, females): <4--5 (normal); >5 suggests hyperandrogenism
  \item Note: SHBG reference ranges vary by assay, laboratory, and patient population. Always use laboratory-specific ranges.
\end{itemize}

\section*{Follow-up Tests}
\begin{itemize}
  \item Total testosterone: Must be measured simultaneously with SHBG. Calculated free and bioavailable testosterone are derived from these values.
  \item Free testosterone (equilibrium dialysis or calculated): The biologically active fraction. Calculation uses SHBG, albumin, and total testosterone in mass action equations.
  \item Fasting glucose, insulin, HOMA-IR: Assess insulin resistance when low SHBG is found. SHBG is a surrogate marker for insulin resistance.
  \item TSH and free T4: Exclude thyroid dysfunction causing SHBG abnormalities.
  \item LH and FSH: Classify gonadal disorders. SHBG helps explain discrepancies between total testosterone and gonadotropin levels.
  \item Estradiol: Elevated SHBG from estrogen therapy will also bind estradiol; total estradiol may be elevated while free estradiol is unchanged.
  \item DHEA-S: In the evaluation of hyperandrogenism with abnormal SHBG and testosterone.
  \item Liver function tests (ALT, AST, albumin): Evaluate for liver disease as a cause of abnormal SHBG.
  \item Lipid profile: Assess metabolic syndrome components associated with low SHBG.
  \item Hemoglobin A1c: Screen for diabetes in patients with low SHBG.
\end{itemize}

\section*{References}
\begin{enumerate}
  \item Bhasin S, Brito JP, Cunningham GR, et al. Testosterone therapy in men with hypogonadism: an Endocrine Society clinical practice guideline. J Clin Endocrinol Metab. 2018;103(5):1715-1744.
  \item Rosner W, Auchus RJ, Azziz R, et al. Utility, limitations, and pitfalls in measuring testosterone: an Endocrine Society position statement. J Clin Endocrinol Metab. 2007;92(2):405-413.
  \item Vermeulen A, Verdonck L, Kaufman JM. A critical evaluation of simple methods for the estimation of free testosterone in serum. J Clin Endocrinol Metab. 1999;84(10):3666-3672.
  \item Teede HJ, Misso ML, Costello MF, et al. International evidence-based guideline for the assessment and management of PCOS. Fertil Steril. 2018;110(3):364-379.
  \item Hammond GL. Diverse roles for sex hormone-binding globulin in reproduction. Biol Reprod. 2011;85(3):431-441.
  \item Le TN, Nestler JE, Strauss JF 3rd, Wickham EP 3rd. Sex hormone-binding globulin and type 2 diabetes mellitus. Trends Endocrinol Metab. 2012;23(1):32-40.
  \item Burtis CA, Ashwood ER, Bruns DE. Tietz Textbook of Clinical Chemistry and Molecular Diagnostics. 6th ed. Elsevier; 2023.
\end{enumerate}

\clearpage
